% ============================================================
%  Chapter 1 -- Introduction
% ============================================================
\chapter{Introduction}
\label{chap:introduction}

% ------------------------------------------------------------
\section{Background: The Automation Paradigm Shift}
\label{sec:intro-background}
% ------------------------------------------------------------

For decades, the dominant fear about automation centered on robots replacing
factory workers and routine physical labor.  The image of a mechanical arm
welding car frames on an assembly line came to represent technological
displacement.  When economists and computer scientists studied automation
risk, they consistently found that blue-collar, repetitive, and physically
demanding jobs were the most vulnerable.  This framing shaped public policy,
workforce development programs, and popular anxiety for an entire generation.

\citet{Frey_Osborne_2017} formalized this understanding in what became the
most cited paper in automation economics.  They hand-labeled 70~occupations
as automatable or not, then trained a Gaussian Process Classifier on O*NET
occupational variables to predict automation probabilities for
702~U.S.\ occupations.  Their central finding: 47\% of total
U.S.\ employment fell into the ``high risk'' category, with automation
probabilities above 70\%.  The jobs at greatest risk were those heavy in
physical manipulation, routine data processing, and structured environments.
They identified three ``engineering bottlenecks'' that protected occupations
from automation: perception and manipulation (fine motor skills in
unstructured environments), creative intelligence (producing novel ideas),
and social intelligence (reading human emotions and negotiating).  In 2013,
these bottlenecks seemed robust.

Between 2022 and 2025, those bottlenecks began to crack.  Large language
models like GPT-4 started producing legal briefs, writing software,
summarizing medical literature, and generating marketing copy that was often
indistinguishable from human-authored work.  GitHub Copilot accelerated
developer productivity by 55\% in controlled experiments
\citep{Peng_Kalliamvakou_2023}.  In a randomized trial at MIT, ChatGPT
users completed mid-level professional writing tasks 37\% faster while
producing higher-quality output \citep{Noy_Zhang_2023}.  The automation
crosshairs shifted from the factory floor to the corporate office.  The jobs
that \citet{Frey_Osborne_2017} had identified as ``safe'' (analytical,
communicative, and cognitive roles) were suddenly the ones most exposed to AI
capability.

The metrics used to discuss automation risk are outdated.  Frey and Osborne's probabilities
reflect 2013-era technology, a world before transformer architectures,
before ChatGPT, and before AI could write poetry or pass the bar exam.
Policymakers and workers still reference their findings, but the
occupational landscape of risk has fundamentally inverted.  I needed a new
metric, one that acknowledges the historical trajectory of automation while
accurately measuring the present capability overhang of Generative~AI.

In 2020, GPT-3
demonstrated that scaling language models to 175~billion parameters produced
emergent capabilities in writing, translation, and reasoning that
researchers had not anticipated at smaller scales.  By November~2022,
ChatGPT attracted 100~million users within two months, the fastest consumer
technology adoption in history.  Within 18~months, GPT-4 passed the Uniform
Bar Examination in the 90th percentile and scored competitively on Advanced
Placement exams in biology, chemistry, and environmental science
\citep{OpenAI_2023}.  Google, Meta, and Anthropic followed with comparable
systems, confirming that the frontier was not a single company's achievement
but a structural property of large transformer architectures trained on
internet-scale data.

The geographic scope of this dissertation covers the United States labor
market, reflecting the data foundations I rely on: O*NET is a
U.S.\ Department of Labor product covering U.S.\ occupational
classifications, and \citet{Frey_Osborne_2017}'s original study analyzed
U.S.\ employment data exclusively.  The structural insights about cognitive
versus manual exposure may generalize to other high-income economies with
similar occupational distributions, but I do not test that claim here.  The
temporal scope spans from the 2013 baseline established by
\citet{Frey_Osborne_2017} to a 2026 counterfactual scenario representing
current GenAI capabilities.  I do not attempt to forecast capabilities
beyond 2026, because the pace of AI development makes such projections
unreliable.

In the broader literature, economists describe a ``Third Wave'' of
automation anxiety.  The First Wave (1960s--1970s) centered on mainframe
computers displacing clerical workers.  The Second Wave (1990s--2010s)
focused on industrial robots and outsourcing threatening manufacturing
employment.  The Third Wave, beginning around 2022, targets the professional
class (lawyers, accountants, programmers, writers, and analysts) whose
educational credentials were previously considered reliable insulation
against technological displacement.  The MAEI is designed to quantify this
Third Wave with the same empirical rigor that \citet{Frey_Osborne_2017}
applied to the Second.

% ------------------------------------------------------------
\section{Research Rationale, Scope, and Research Contributions}
\label{sec:intro-rationale}
% ------------------------------------------------------------

A gap exists in the existing literature.  No published work
combines all four of the following elements: (1)~an updated automation risk
framework recalibrated from Frey and Osborne's original structure for modern
AI capabilities, (2)~natural language processing on occupational task
descriptions to extract semantic features, (3)~scenario-based capability
multipliers with Monte Carlo uncertainty quantification, and (4)~external
validation against an independently constructed AI exposure benchmark.
Several papers address one or two of these elements in isolation.
\citet{Felten_Raj_Seamans_2021} constructed an AI exposure index but used
crowdsourced linkages rather than scenario modeling.
\citet{Eloundou_Manning_2023} estimated LLM exposure breadth but did not
build a reproducible occupation-level index.  \citet{Webb_2020} used NLP
patent-matching but did not anchor to historical baselines.
\citet{Anthropic_2026} measured observed AI usage patterns but relied on
proprietary data that cannot be replicated academically.  The MAEI fills
this gap by integrating all four elements into a single reproducible
pipeline using only publicly available data.

This dissertation makes three specific contributions to the field:

\textbf{First}, I construct the MAEI itself: a reproducible structural index
covering 1,016~U.S.\ occupations, built entirely from public O*NET data and
validated against an independent benchmark.

\textbf{Second}, I introduce a three-layer scoring architecture that
honestly separates what the machine learning model contributes (structural
validation) from what the scenario logic contributes (occupation rankings).
This decomposition emerged from discovering that the ensemble's tree-based
learners saturate on multiplied features, and rather than hiding this
finding, I designed the architecture to acknowledge and account for it.

\textbf{Third}, I provide empirical proof of structural stability through a
seven-scenario multiplier ablation study demonstrating that the occupation
rankings are invariant to specific multiplier values ($\rho$ range:
0.7636--0.7651 across all scenarios).  This converts a potential weakness
(subjective multiplier choice) into a demonstrated strength (structural
stability).

% ------------------------------------------------------------
\section{Research Statement and Research Questions}
\label{sec:intro-rqs}
% ------------------------------------------------------------

\textbf{Central Thesis:} The MAEI quantifies occupational overlap with
Generative AI capabilities and shows that the
occupational profile of automation risk has inverted since 2013, shifting
from routine physical labor to cognitive, analytical, and communicative work.

Three research questions guide this investigation:

\textbf{RQ1:} Can a machine learning model trained on 2013 automation labels
and modern O*NET features produce a structurally plausible baseline for
measuring AI exposure?

\textbf{RQ2:} Does a scenario-based perturbation of cognitive and analytical
features generate an occupation ranking that converges with independently
constructed AI exposure benchmarks?

\textbf{RQ3:} How robust are the resulting rankings to changes in the
underlying multiplier assumptions?

% ------------------------------------------------------------
\section{Scope and Clarification}
\label{sec:intro-scope}
% ------------------------------------------------------------

Before proceeding, I want to clarify what the MAEI is and what it is not.
The MAEI measures capability overlap, the degree to which an occupation's
daily tasks map onto the current capabilities of Generative~AI.  A high MAEI
score for ``Lawyer'' means that AI systems can now perform many tasks lawyers
do daily (document review, case synthesis, contract drafting).  It does not
mean that lawyers will be unemployed.  As \citet{Acemoglu_Restrepo_2018}
have shown, the relationship between task automation and actual job
displacement is mediated by economic adoption costs, regulatory friction,
organizational inertia, and worker adaptability.  Anthropic's
\citeyearpar{Anthropic_2026} recent empirical analysis found ``no systematic
unemployment increase'' for workers in highly AI-exposed occupations, even
though capability overlap is substantial.  The MAEI is a scenario-based
structural index, not an econometric forecast of unemployment.

% ------------------------------------------------------------
\section{Chapter Outline}
\label{sec:intro-outline}
% ------------------------------------------------------------

Chapter~2 reviews the theoretical and empirical literature underpinning this
project, from \citet{Frey_Osborne_2017}'s foundational work through modern
GenAI exposure research.  Chapter~3 describes the research methodology,
including data engineering, NLP feature extraction, ensemble modeling, the
three-layer MAEI architecture, and Monte Carlo uncertainty quantification.
Chapter~4 presents the results, including baseline model performance,
external validation, multiplier ablation, and the tree saturation analysis.
Chapter~5 discusses limitations, proposes future research directions, and
concludes.
